\exercisegroup

\begin{enumerate}
\def\labelenumi{\arabic{enumi})}
\tightlist
\item
   Let K be one of the topological fields R, C, H. Let V be a K-vector space possessing a countable generating system, equipped with its finest topology as a topological vector space (EVT, I, §1, p. 11, cor. 4). For every integer \(n \geqslant 1\) , let \(\mathrm{B}_{n}(\mathrm{V})\) be the subspace of \(V^{n}\) consisting of free sequences and \(\mathrm{G}_{n}(\mathrm{V})\) the set of subspaces of dimension n of V. Denote by \(p: \mathrm{B}_{n}(\mathrm{V}) \to \mathrm{G}_{n}(\mathrm{V})\) the map which associates to a sequence \(v \in \mathrm{B}_{n}(\mathrm{V})\) the vector subspace it generates, and equip the set \(\mathrm{G}_{n}(\mathrm{V})\) with the coarsest topology for which the map p is continuous.
\end{enumerate}

\begin{enumerate}
\def\labelenumi{\alph{enumi})}
\item
   Prove that one defines a right action of the group \(\mathbf{GL}(n,\mathrm{K})\) on \(\mathrm{B}_{n}(\mathrm{V})\) by setting \(v \cdot g = (\sum_{i=1}^{n} v_{i} a_{i,j})_{1 \leqslant j \leqslant n}\) for \(v = (v_{1}, \ldots, v_{n}) \in \mathrm{B}_{n}(\mathrm{V})\) and \(g = (a_{i,j}) \in \mathbf{GL}(n, \mathrm{K})\) .
\item
  Prove that the map p induces a homeomorphism of the quotient space \(\mathrm{B}_{n}(\mathrm{V})/\mathbf{G}\mathbf{L}(n,\mathrm{K})\) onto \(\mathrm{G}_{n}(\mathrm{V})\) . Deduce that the space \(\mathrm{G}_{n}(\mathrm{V})\) is paracompact, and metrizable if V is of finite dimension.
\item
  Prove that \(\mathrm{B}_{n}(\mathrm{V})\) is a principal fiber bundle with group \(\mathbf{GL}(n,\mathrm{K})\) and base \(\mathbf{G}_{n}(\mathrm{V})\) .
\item
  Assume that V is of infinite dimension. Prove that the space \(\mathrm{B}_{n}(\mathrm{V})\) is contractible. Deduce that the space \(\mathrm{G}_{n}(\mathrm{V})\) is a classifying space for the group \(\mathbf{GL}(n,\mathrm{K})\) .
\end{enumerate}

\begin{enumerate}
\def\labelenumi{\arabic{enumi})}
\setcounter{enumi}{1}
\tightlist
\item
  Let H be a separated pre-Hilbert space of infinite dimension. Let \((e_{i})_{i\geqslant0}\) be an orthonormal family of vectors of H which generates a dense subspace of H. For every integer \(n\geqslant1\) , let \(V_{n}(H)\) be the subspace of \(H^{n}\) consisting of sequences \((u_{1},\ldots,u_{n})\) which are orthonormal.
\end{enumerate}

\begin{enumerate}
\def\labelenumi{\alph{enumi})}
\item
  Prove that there exists a unique continuous linear map T: H → H such that \(\mathrm{T}(e_{i}) = e_{i+1}\) for every integer i. Prove that one has \(\|T(x)\| = \|x\|\) for all \(x \in H\) .
\item
  Let H' be the subspace of H formed by the vectors orthogonal to \(e_{i}\) , for i \textless{} n. Prove that the following conditions are equivalent: (i) x and \(\mathrm{T}^{n}(x)\) are linearly dependent; (ii) \(\mathrm{T}(x) = x\) ; (iii) \(x \in H'\) .
\item
  Prove that there exists a unique homotopy \(\sigma\colon\mathrm{V}_{n}(\mathrm{H})\times\mathbf{I}\to\mathrm{V}_{n}(\mathrm{H})\) such that, for every \((u_{1},\ldots,u_{n})\in\mathrm{V}_{n}(\mathrm{H})\) and every \(t\in\mathbf{I}\) , \(\sigma((u_{1},\ldots,u_{n}),t)\) is the sequence deduced from the family \((1-t)(u_{1},\ldots,u_{n})+t(\mathrm{T}^{n}(u_{1}),\ldots,\mathrm{T}^{n}(u_{n}))\) by the orthonormalization process (EVT, V, p. 23, prop. 6).
\item
  Prove that there exists a unique homotopy \(\tau\colon\mathrm{V}_{n}(\mathrm{H}^{\prime})\times\mathbf{I}\to\mathrm{V}_{n}(\mathrm{H})\) such that, for every \((u_{1},\ldots,u_{n})\in\mathrm{V}_{n}(\mathrm{H}^{\prime})\) and every \(t\in\mathbf{I}\) , \(\sigma((u_{1},\ldots,u_{n}),t)\) is the sequence deduced from the family \((1-t)(u_{1},\ldots,u_{n})+t(e_{1},\ldots,e_{n})\) by the orthonormalization process.
\item
  Prove that the space \(V_{n}(H)\) is contractible.
\end{enumerate}

\begin{enumerate}
\def\labelenumi{\arabic{enumi})}
\setcounter{enumi}{2}
\tightlist
\item
  Let H be a separated pre-Hilbert space. Let n be an integer \(\geqslant 1\) ; we denote by \(V_{n}(H)\) the subspace of \(H^{n}\) consisting of orthonormal sequences. Let G be a closed subgroup of \(O(n, \mathbf{R})\) .
\end{enumerate}

\begin{enumerate}
\def\labelenumi{\alph{enumi})}
\tightlist
\item
   Prove that the orthogonal group O(n, R) (LIE, IX, §3, n° 5, p. 22) acts properly and freely on the space \(V_{n}(H)\) by the formula
\end{enumerate}

\[
(u _ {1}, \dots , u _ {n}) \cdot (a _ {i, j}) = \left(\sum_ {i = 1} ^ {n} a _ {i, j} u _ {i}\right).
\]

\begin{enumerate}
\def\labelenumi{\alph{enumi})}
\setcounter{enumi}{1}
\item
  Let G be a closed subgroup of \(O(n,\mathbf{R})\) . Prove that the canonical surjection makes \(V_{n}(\mathrm{H})\) a principal fiber bundle with group G and base \(V_{n}(\mathrm{H})/\mathrm{G}\) (Begin by treating the case where \(G = O(n,\mathbf{R})\) .)
\item
  Prove that the quotient space \(V_{n}(H)/G\) is metrizable. (Endow \(H^{n}\) with its natural structure of pre-Hilbert space and consider the gauge d on \(V_{n}(H)\) given by \(d(u,v)=\inf_{g\in G}\|u-v\cdot g\|\) .)
\item
  Prove that the quotient space \(V_{n}(H)/G\) is a classifying space for G.
\item
  Deduce that every compact Lie group has a classifying space which is a metrizable topological space. (Apply LIE, IX, §9, no. 2, p. 91, th. 1.)
\end{enumerate}

\begin{enumerate}
\def\labelenumi{\arabic{enumi})}
\setcounter{enumi}{3}
\item
  By considering a complex Hilbert space, construct in an analogous manner a metrizable classifying space for the unitary group U(n, C) (LIE, IX, §, no. 4, p. 21)
\item
  \begin{enumerate}
  \def\labelenumii{\alph{enumii})}
  \tightlist
  \item
    Let X be a paracompact topological space. Prove that every principal fiber bundle with base X and group R is trivializable.
  \end{enumerate}
\end{enumerate}

\begin{enumerate}
\def\labelenumi{\alph{enumi})}
\setcounter{enumi}{1}
\tightlist
\item
  Let X be the Alexandroff half-line (TG, IV, p. 49, exerc. 12). Construct a principal fiber bundle with base X and group R which is not trivializable.
\end{enumerate}

\begin{enumerate}
\def\labelenumi{\arabic{enumi})}
\setcounter{enumi}{5}
\tightlist
\item
  Let X be the quotient space of the space \(R \times \{0,1\}\) by the finest equivalence relation for which the points \((t,0)\) and \((t,1)\) are equivalent, for every \(x \in R_{+}^{*}\) . We denote by p: \(R \times \{0,1\} \to X\) the canonical surjection.
\end{enumerate}

\begin{enumerate}
\def\labelenumi{\alph{enumi})}
\item
  Prove that the topological space X is homotopic to a point.
\item
  For which points \(x \in X\) is the pointed space (X, x) contractible?
\item
  Let \(U = p(\mathbf{R} \times \{0\})\) and \(V = p(\mathbf{R} \times \{\mathrm{V}\})\) . Let G be a topological group. Construct a bijection from the set \(\mathcal{C}(\mathbf{R}_{+}^{*}; \mathbf{G})\) of continuous maps from \(R_{+}^{*}\) into G onto the set of isomorphism classes of principal fiber bundles with base X and group G.
\item
  Deduce that there exist principal fiber bundles with base X and group R which are not trivializable.
\end{enumerate}

\begin{enumerate}
\def\labelenumi{\arabic{enumi})}
\setcounter{enumi}{6}
\tightlist
\item
  Let G be a topological group and let E be a principal fiber bundle with base \(S_{2}\) and group G. We denote by \(S_{2}^{+}\) and \(S_{2}^{-}\) the two hemispheres, intersections of \(S_{2}\) and the half-spaces defined by the conditions \(z \geqslant 0\) and \(z \leqslant 0\) respectively; we identify \(S_{2}^{+} \cap S_{2}^{-}\) with \(S_{1}\) .
\end{enumerate}

\begin{enumerate}
\def\labelenumi{\alph{enumi})}
\item
  Prove that the bundle E has a section \(s^{+}\) (resp. \(s^{-}\) ) over \(S_{2}^{+}\) (resp. \(S_{2}^{-}\) ). Prove that there exists a unique continuous map \(f\colon S_{1}\to G\) such that \(s^{+}(x)=s^{-}(x)\cdot f(x)\) for all \(x\in S_{1}\) .
\item
  Prove that the strict homotopy class of the map from \(S_{1}\) into G given by \(x \mapsto c(x)c(1)^{-1}\) does not depend on the choice of sections \(s^{+}\) and \(s^{-}\) ; it is denoted \(c(\mathrm{E})\) .
\item
  Prove that two principal fiber bundles E and \(E'\) , with base \(S_{2}\) and group G, are isomorphic if and only if \(c(\mathrm{E}) = c(\mathrm{E}')\) .
\item
  Prove that, for every class \(c \in \pi_{1}(G, e)\) , there exists a principal fiber bundle E, with base \(S_{2}\) and group G such that \(c(E) = c\) .
\end{enumerate}

\begin{enumerate}
\def\labelenumi{\arabic{enumi})}
\setcounter{enumi}{7}
\tightlist
\item
  Let X be a paracompact topological space and let S(X) be its suspension (I, p. 149, exerc. 1); we denote by \(p\colon X \times [0,1] \to S(X)\) the canonical surjection. Let \(a\) be a point of X.
\end{enumerate}

\begin{enumerate}
\def\labelenumi{\alph{enumi})}
\item
  Let G be a topological group and let E be a principal fiber bundle with group G and base S(X). Prove that the restriction of E to the subspace \(p(\mathrm{X} \times [0,1/2])\) admits a continuous section \(s_{0}\) , and that the restriction of E to the subspace \(p(\mathrm{X} \times [1/2,1])\) admits a continuous section \(s_{1}\) .
\item
  Prove that there exists a unique map \(f\colon \mathrm{X}\to \mathrm{G}\) such that \(s_1(x) = s_0(x)\cdot f(x)\) for all \(x\in \mathrm{X}\) .
\item
  Prove that the strict homotopy class in \([(X,a);(G,e)]\) of the pointed map \(x \mapsto f(x)f(a)^{-1}\) does not depend on the choice of sections \(s_{0}\) and \(s_{1}\) ; it is denoted \(c(\mathrm{E})\) .
\item
  Prove that two principal fiber bundles E and \(E'\) , with base S(X) and group G, are isomorphic if and only if \(c(\mathrm{E}) = c(\mathrm{E}')\) .
\item
  Prove that, for every class \(c \in [(X, a); (G, e)]\) , there exists a principal fiber bundle E, with base S(X) and group G, such that \(c(E) = c\) .
\end{enumerate}

\begin{enumerate}
\def\labelenumi{\arabic{enumi})}
\setcounter{enumi}{8}
\tightlist
\item
  We identify the complex projective line \(\mathbf{P}^{1}(\mathbf{C})\) with the set of lines of \(C^{2}\) . Let E be the subspace of \(\mathbf{C}^{2} \times \mathbf{P}^{1}(\mathbf{C})\) consisting of the pairs \(((u, v), x)\) such that \(v \in x\) and \(|u|^{2} + |v|^{2} = 1\) . We equip E with the action of the group \(S_{1}\) given by \(((u, v), x) \cdot z = ((zu, zv), x)\) .
\end{enumerate}

\begin{enumerate}
\def\labelenumi{\alph{enumi})}
\item
  Prove that the second projection \(\mathrm{pr}_{2}\colon\mathrm{E}\to\mathbf{P}^{1}(\mathbf{C})\) endows E with a structure of principal fiber bundle with base \(\mathbf{P}^{1}(\mathbf{C})\) and group \(S_{1}\) .
\item
  Let \(\varphi\) be a homeomorphism of \(\mathbf{S}_2\) onto \(\mathbf{P}^1 (\mathbf{C})\) ; prove that the group \(\pi_1(\mathbf{S}_1,e)\) is generated by the class \(c(\varphi^{*}\mathrm{E})\) .
\end{enumerate}

\backmatter

